\section{Problem analysis}

For the player to have as little interaction with the robot as possible, it is therefore required to detect the player's moves and keep track of where the chessboard is located in reference to the robot, in case the player moves the chessboard during a move. To solve these 2 problems, a camera is used to detect the move the player has made, as well as the location of the chessboard. The pieces on a chessboard are usually placed close to each other, the gripper therefore requires both high precision and high accuracy to prevent knocking down other pieces. When the robot captures an opponent's piece, it will also have to account for human placement inaccuracy. This means the gripper must be able to pick up a piece regardless of its position on a tile.

For moving the pieces a Universal Robots' UR5 cb3 is used.\cite{UR5} The robot can be controlled using various means, the most popular of which is the teaching pendant. Using this, pre-planned movements between known points can be written in an intuitive language in programs called UR programs. However, as the robot needs to move dynamically between squares and to account for the dynamic placement of the chessboard, simple programs onboard the teaching pendant are not enough. Hence ROS2 is used as it handles communication, inverse kinematics, and motionplanning for the robot. Furthermore ROS2 provides a graphical user interface assisting in development.

To control the DC motor, a Raspberry Pi Pico \cite{Rasperry-Pi-Pico} is used due to its GPIO (General Purpose Input Output) pins to control the motor, as well as the team's prior experience with the Pico. When using a single DC motor for opening and closing the gripper, an H-bridge can be used to produce a bidirectional current.

For the robot to apply the best move, it needs to interact with an engine like Stockfish \cite{Stockfish}, which dictates the robot's move. Therefore, the robot is required to keep track of the position of all pieces in order to communicate with the engine.
\begin{table}[b]
    \centering
    \begin{tabular}{p{13cm}|c}
        \toprule
        \rowcolor{gray!30}
        \textbf{Technical Requirements} & \textbf{Success rate} \\
        \midrule

        \rowcolor{gray!15}
        \textbf{Gripper} &\\
        The gripper has to be able to pick up any piece, without knocking down other pieces & 98\%\\
        The gripper has to be able to let any piece go at the designated square & 98\%\\
        The gripper has to know what state it is in (open, closed or on standby) & 100\% \\
  
        
        \rowcolor{gray!15}
        \textbf{Camera} &\\
        The camera has to identify where the board is located. & 98\% \\
        The camera has to identify which piece has been moved since the last frame & 98\% \\
        
        \rowcolor{gray!15}
        \textbf{Robot} &\\
        The robot movement must not knock down other pieces & 100\% \\
        
        \bottomrule
    \end{tabular}
    \caption{Technical Requirements for the project}
    \label{tab:project_requirements}
\end{table}

\subsection{Chess and piece specifications}
\label{sec:ChessSpecs}
To design the gripper, the specifications of the chess pieces and the tile dimensions are measured to ensure the gripper does not knock down surrounding pieces. 
The height of the chess pieces varies from 3 cm to 5.3 cm, with the smallest being a pawn and the tallest being the king. The diameter of the bottom varies between 1.6 cm and 1.9 cm. Furthermore, every piece have a small indent in the bottom of the piece in the same height.
The tiles' width and length are measured to be 3.6 cm, and therefore making the width and length of the board $3.6 \space \text{cm} \cdot 8 = 28.8 \text{cm}$

\subsection{Technical requirements}
For testing purposes, technical requirements for measuring and determining the success of the project are created. A game of chess is usually around 50 moves, although this varies. The desired success rate is determined with this in mind.
These can be seen in \cref{tab:project_requirements}.






